\subsection{El coeficiente del perfil y su estabilidad}
Para sumar las dos correcciones de G1 conviene expresarlas en el mismo perfil $\nu=kP$.
Al componer núcleos usamos las medidas uniformes de filas y columnas, de modo que cada suma
sobre una fibra queda dividida por el tamaño del tablero. Escriba $B=\nu-1$, $A_R=BB^\top$, $A_C=B^\top B$,
$T_X=A_X(I-A_X)^{-1}$, $U=T_RB=BT_C$ y $V=B+U$. Para $X=R,C$ sea
$j_X(t)=(T_X-A_X)(t,t)$, $t_X(t)=T_X(t,t)$ y $j_0=\int j_R=\int j_C$.
Los cubos en integrales dobles son potencias de los valores de los núcleos. En esta notación,
el funcional que entra en (T) es
\begin{align*}
 \mathfrak F={}&\tfrac12j_0^2+\operatorname{tr}(T_R^2)-\tfrac32 j_0
    -\tfrac34\sum_X\int j_X^2+\langle j_R,V j_C\rangle\\
 &-\tfrac12\sum_X\langle j_X,T_Xj_X\rangle-\tfrac13\sum_X\iint T_X^3
    +2\iint BU^2+\tfrac23\iint U^3,\qquad
 \Phi=\mathfrak F-\tfrac1{12}.
\end{align*}
Sea $\bar e$ la proyección de $\nu$ en $L^2(\nu)$ sobre el ortogonal de las funciones
aditivas de fila y columna. Ponga $m_R=\int\nu^2\,dy$, $m_C=\int\nu^2\,dx$. La parte
determinista del coeficiente es
\begin{align*}
 \delta_1={}&\tfrac1{12}+\tfrac56\iint\nu^3+\tfrac12\iint\nu\bar e^2
 -\iint V\nu\bar e-\tfrac12\iint\nu^2\\
 &+\tfrac12\iint\nu^2(t_R+t_C)-\iint\nu^2V-\tfrac14\sum_X\int m_X^2 .
\end{align*}
Para obtener $\delta_1$ de la normalización, relacionamos el punto estacionario $p$ con el
escalado $P$. Existen potenciales $\alpha,\beta$ con
$P_{ij}=\exp(\alpha_i+\beta_j)p_{ij}(1-p_{ij})$. Multiplicar el logaritmo por $P_{ij}$
y sumar elimina el calibre. Sustituyendo en la expresión exacta de $\Delta$ y expandiendo
la entropía y el pseudodeterminante se obtiene $\Delta=\delta_1/k+O(k^{-2})$.
La expansión emplea la difusión para acotar las sumas de potencias y la brecha para
acotar las inversas. Para verificar el coeficiente, úsense $V B^\top=T_R$, $B^\top V=T_C$,
$T_X A_X=T_X-A_X$ y que $B$ tiene marginales cero. Si $g=\iint B^2$ y
$\gamma_X(t)=A_X(t,t)$, la proyección anterior se escribe
\[
 \bar e=B+g-\rho_R-\rho_C,\quad
 \rho_R=(I+T_R)\gamma_R-V\gamma_C,\quad
 \rho_C=(I+T_C)\gamma_C-V^\top\gamma_R .
\]
Sustituir esta expresión y $t_X=\gamma_X+j_X$ en $\delta_1$ y en la fórmula estática da
$\mathfrak f_{\rm e}+\delta_1=\Phi$ en cada malla. Por tanto
$R=R_{\rm est}+R_{\rm det}$, con $R_{\rm est}=\mathcal E(P)-\mathfrak f_{\rm e}/k$ y
$R_{\rm det}=\Delta-\delta_1/k$. Ambos son $O(k^{-2})$, lo que prueba (T). En el cociente de
dos permanentes aparece la diferencia de sus errores. La descomposición permite estimar esa
diferencia por separado para $R_{\rm est}$ y $R_{\rm det}$.

\paragraph{La estabilidad estática $(L^{*+})$ [D].}
En $(L^*)$ bastaba acotar todos los diagramas de exceso al menos dos. Una vez cambiado el
tablero, hay que reconocer cuáles de sus contribuciones sobreviven en el menor. Consérvense
los diagramas de excesos dos y tres, con la corrección de índices coincidentes
$-\frac56\sum C_{aa}^3$. Los mismos mayorantes dejan una cola $O(k^{-4})$.
El promedio gamma y el logaritmo dan $R_{\rm est}=f_2/k^2+f_3/k^3+O(k^{-4})$, con
$f_2,f_3$ uniformemente acotados. Para pasar a un menor, el defecto de masa por fibra es
$O(k^{-1})$. Los potenciales que lo corrigen tienen norma supremo $O(k^{-1})$.
Al quitar una sola casilla el defecto se concentra en dos fibras. La inversa en el calibre tiene
norma supremo acotada y entradas fuera de la diagonal $O(k^{-1})$. Los potenciales son
$O(k^{-1})$ en las dos fibras y $O(k^{-2})$ en el resto. El argumento de perturbación conserva
la clase con cotas $2M_0,\gamma/2$. Al expandir la diferencia de un producto de propagadores
se marca la arista cambiada. Las sumas restantes se acotan sobre un árbol generador, usando
la norma absoluta de fila y la cota de entrada. La restricción del índice a la fibra afectada
ahorra un factor $k^{-1}$ en el borrado localizado. Los coeficientes $f_2,f_3$ cambian en
$O(k^{-1})$ bajo el menor y en $O(k^{-2})$ bajo el borrado. Al variar también los denominadores,
esto da los órdenes $k^{-3}$ y $k^{-4}$ de los respectivos incrementos de $R_{\rm est}$.
